\uuid{9avh}
\titre{ Rétropropagation : ordre des opérations }
\niveau{L3}
\module{Analyse de données}
\chapitre{Réseaux de neurones}
\sousChapitre{Rétropropagation du gradient}
\theme{Backpropagation, ordre d'exécution, erreur de variable non définie}
\auteur{Maxime Nguyen}
\datecreate{2026-03-19}
\organisation{AMSCC}
\difficulte{2}

\contenu{
\texte{
  On compare deux implémentations du backward pass d'un MLP à une couche cachée :
  \begin{itemize}
    \item[A.] \begin{verbatim}
dW2 = (1/m) * A1.T @ dy
dA1 = dy @ W2.T
dZ1 = relu_backward(dA1, Z1)
dW1 = (1/m) * X.T @ dZ1
\end{verbatim}
    \item[B.] \begin{verbatim}
dW1 = (1/m) * X.T @ dZ1
dZ1 = relu_backward(dA1, Z1)
dA1 = dy @ W2.T
dW2 = (1/m) * A1.T @ dy
\end{verbatim}
  \end{itemize}
}
\begin{enumerate}

\item
  \question{Lequel peut fonctionner ? Lequel produira une erreur à l'exécution, et pourquoi ?}
  \reponse{
    \textbf{A} est correct : les variables sont calculées dans le bon ordre — \texttt{dA1} est défini avant
    d'être utilisé dans \texttt{relu\_backward}, et \texttt{dZ1} est défini avant \texttt{dW1}.

    \textbf{B} produira une erreur \texttt{NameError} (ou utilisera des valeurs périmées) :
    \texttt{dW1} est calculé en utilisant \texttt{dZ1} avant que \texttt{dZ1} ne soit défini,
    et \texttt{dZ1} utilise \texttt{dA1} avant que \texttt{dA1} ne soit défini.
    La rétropropagation doit impérativement se faire dans l'ordre inverse du forward pass.
  }

\end{enumerate}
}
